\documentclass[10pt]{article}
\usepackage[margin=1in, paperwidth=8.5in, paperheight=11in]{geometry}
\usepackage{ifpdf,amsmath, amssymb, comment, color, graphicx, stmaryrd,setspace,enumitem,tikz, fancyhdr, wrapfig, textcomp, units, mathptmx, pdfpages, cancel, multicol}
\usepackage{hyperref}
\hypersetup{
    colorlinks=true,
    urlcolor=blue,
}

\setlength{\headheight}{14.5pt}
\newcommand{\Q}{\mathbb{Q}}
\newcommand{\R}{\mathbb{R}}
\newcommand{\Z}{\mathbb{Z}}
\newcommand{\vu}{\mathbf{u}}
\newcommand{\vv}{\mathbf{v}}
\newcommand{\vw}{\mathbf{w}}
\newcommand{\vi}{\mathbf{i}}
\newcommand{\vj}{\mathbf{j}}
\newcommand{\vk}{\mathbf{k}}
\newcommand{\vn}{\mathbf{n}}
\newcommand{\vr}{\mathbf{r}}
\newcommand{\va}{\mathbf{a}}
\newcommand{\vF}{\mathbf{F}}
\newcommand{\vL}{\mathbf{L}}
\newcommand{\vT}{\mathbf{T}}
\newcommand{\vN}{\mathbf{N}}
\newcommand{\proj}{\operatorname{proj}}
\newcommand{\orth}{\operatorname{orth}}
\newcommand\del\nabla
\newcommand\dotp[1][.5]{\,\mathbin{\vcenter{\hbox{\scalebox{#1}{$\bullet$}}}}\,}
\newcommand{\slice}{_\text{slice}}

% Solution text is in red. If you want the solutions to show, remove the \iffalse from the definition of the \red command.
%\newcommand{\red}[1]{ %\iffalse
%	\textcolor{red}{#1} }%\fi}

\newenvironment{red}{\color{red}}{\ignorespacesafterend}
\newenvironment{blue}{\color{blue}}{\ignorespacesafterend}
\newcommand{\green}[1]{\textcolor{green}{#1}}
\renewcommand{\section}[1]{\begin{center} \textbf{#1} \\\end{center}}
%
\hyphenpenalty=5000
\setlength{\parindent}{0in}
%\oddsidemargin=-.25in
\allowdisplaybreaks
\pagestyle{fancy}
\renewcommand{\headrulewidth}{0pt}
\lhead{MATH 202}
\rhead{Spring 2024}
%\lfoot{\copyright\ CLEAR Calculus 2010}
\cfoot{}

\begin{document}
%


%\onehalfspacing
\allowdisplaybreaks
%##################################################################
\section{PS\#8: Improper integrals - \red{Answer key} }

\everymath{\displaystyle}
\begin{enumerate}[leftmargin=0pt]

\item (\href{https://activecalculus.org/single/sec-6-5-improper.html#IZN}{Activity 6.5.4}) For each of the following definite integrals, decide whether the integral is improper or not. If the integral is proper, evaluate it using the First FTC. If the integral is improper, determine whether or not the integral converges or diverges; if the integral converges, find its exact value.
\begin{enumerate}
    \item $\displaystyle \int_0^1 \frac{1}{x^{1/3}} \, dx$
    \begin{red} -- 
        \textbf{improper}, \textbf{converges} to $\dfrac{3}{2}$
    \end{red}

    \item $\displaystyle \int_0^2 e^{-x} \, dx$
    \begin{red} -- 
        \textbf{proper}, \textbf{converges} to $-e^{-2}+1$
    \end{red}

    \item $\displaystyle \int_1^4 \frac{1}{\sqrt{4-x}} \, dx$
    \begin{red} -- 
        \textbf{improper}, \textbf{converges} to $2 \sqrt{3}$
    \end{red}
    
    \item $\displaystyle \int_{-2}^2 \frac{1}{x^2} \, dx$
    \begin{red} -- 
        \textbf{improper}, \textbf{diverges}
    \end{red}
    
    \item $\displaystyle \int_0^{\pi/2} \tan(x) \, dx$
    \begin{red} -- 
        \textbf{improper}, \textbf{diverges}
    \end{red}
    
    \item $\displaystyle \int_0^1 \frac{1}{\sqrt{1-x^2}} \, dx$
    \begin{red} -- 
        \textbf{improper}, \textbf{converges} to $\dfrac\pi2$
    \end{red}
    
\end{enumerate}

\item (Modified from \href{https://activecalculus.org/single/sec-7-1-diff-eq-intro.html#mvy}{AC 7.1 Exercise 6}) A magic seashell has whispered into our ear that the following functions might be solutions to the differential equation. However, it is a \textit{rude} magic seashell which sometimes lies. Determine whether each of the following functions is a solution to the differential equation:


\begin{enumerate}
    \item $y(t) = t-2e^t+1$
    
    \begin{red}
    The LHS of the differential equation is $\dfrac{dy}{dt}$, which in this case is $1-2e^t+0$. 
    
    The RHS of the differential equation is $y-t$, which in this case is $\left(t-2e^t+1\right)-t = 1-2e^t$. 
    
    So, this \textbf{is a solution}, yay, because the LHS and the RHS match.
    \end{red}
    \item $y(t) = t+1$
    
    \begin{red}
    The LHS of the differential equation is $\dfrac{dy}{dt}$, which in this case is $1 + 0$. 
    
    The RHS of the differential equation is $y-t$, which in this case is $\left(t+1\right)-t = 1$. 
    
    So, this \textbf{is a solution}, yay, because the LHS and the RHS match.
    \end{red}

    \item $y(t) = 2e^t$

    \begin{red}
    The LHS of the differential equation is $\dfrac{dy}{dt}$, which in this case is $2e^t$. 
    
    The RHS of the differential equation is $y-t$, which in this case is $\left(2e^t\right)-t$. 
    
    So, this \textbf{is not a solution}, boo, because the LHS and the RHS \textbf{don't} match.
    \end{red}
    \item $y(t) = 1+\frac{e^t}{2}+t$
    
    \begin{red}
    The LHS of the differential equation is $\dfrac{dy}{dt}$, which in this case is $0 + \frac{e^t}{2}+1$. 
    
    The RHS of the differential equation is $y-t$, which in this case is $\left(1+\frac{e^t}{2}+t\right)-t = 1+\frac{e^t}{2}$. 
    
    So, this \textbf{is a solution}, yay, because the LHS and the RHS match.
    \end{red}
\end{enumerate}

\end{enumerate}
\end{document}